    \begin{frame}
        \frametitle{Passage de la tension au nombre aléatoire}

        \small

        \begin{block}{1. Conversion binaire → entier}
        \[
        \begin{aligned}
        \underbrace{\texttt{00111110101110111011011110000000}}_{\;1052489600} \\
        \underbrace{\texttt{11111110110101001101101100000000}}_{\;4275362560} \\
        \end{aligned}
        \]
        \end{block}

        \vspace{0.3cm}

        \begin{block}{2. Normalisation sur $[0,1]$}
        \[
        \begin{aligned}
        \frac{1052489600}{2^{32}} &\approx 0{,}245051831 \\
        \frac{4275362560}{2^{32}} &\approx 0{,}995435417 \\
        \end{aligned}
        \]
        \end{block}
    \end{frame}


    \begin{frame}
        \frametitle{Quelque chiffre}
        \begin{center}
            Resistance = 627 690$\Omega$
            \\
            Bobine = 1H
            \\
            Condensateur = 2.5$\cdot 10^{-8}$F
            \\
            Temps d'enregistrement ~1h
            \\
            6 099 998 point enregistrer
            \\
            762 496 octets de donnée
            \\
            190 624 nombre aléatoire
        \end{center}
    \end{frame}
    $
    \begin{cases}
    u_0 = 55\\
    u_{n+1} = (23 \cdot u_n + 19) \bmod 60006 
    \end{cases}
    $